\chapter*{参考文献}
\addcontentsline{toc}{chapter}{参考文献}
\markboth{参考文献}{参考文献}

以下列出编写中参考的主要教材、讲义及习题参考资料，亦供进一步阅读。

\begin{list}{[\arabic{enumi}]}{%
  \usecounter{enumi}%
  \setlength{\leftmargin}{2.4em}%
  \setlength{\labelwidth}{2em}%
  \setlength{\labelsep}{.4em}%
  \setlength{\itemsep}{6pt}%
  \setlength{\parsep}{0pt}%
  \setlength{\topsep}{6pt}}

\item Abbott, Stephen.
\href{https://doi.org/10.1007/978-1-4939-2712-8}{\emph{Understanding Analysis}}.
2nd ed. Springer, 2015.

\item Axler, Sheldon.
\href{https://linear.axler.net/LADR4e.pdf}{\emph{Linear Algebra Done Right}}.
4th ed. Springer, 2024. 参阅作者2026年8月16日更新的电子版。

\item Bobenko, Alexander; R\"orig, Thilo; Springborn, Boris.
\emph{Geometry I}. TU Berlin课程讲义，2021年3月17日稿。

\item Artin, Michael.
\href{https://math.mit.edu/classes/18.721/ag-jun24-2021.pdf}{\emph{Notes for a Course in Algebraic Geometry}}.
MIT 18.721课程讲义，2021年6月24日稿。

\item Grinstead, Charles M.; Snell, J. Laurie.
\href{https://math.dartmouth.edu/~prob/prob/prob.pdf}{\emph{Grinstead and Snell's Introduction to Probability}}.
The CHANCE Project，2006年7月4日修订版。

\item 闻杰．《神奇的圆锥曲线与解题秘诀》．第2版．浙江大学出版社，2023．

\item 张永辉主编．《高考数学培优40讲：解析几何》．清华大学出版社，2023．

\item 张永辉主编．《高考数学培优40讲：函数与导数》．清华大学出版社，2022．

\item 吕荣春主编．《高观点下函数导数压轴题的系统性解读》．电子科技大学出版社，2017．

\item 张永辉主编．《全国卷满分秘籍：导数篇》．清华大学出版社，2017．

\end{list}
